\documentclass[12pt]{article}
\usepackage[T1]{fontenc}
\usepackage[utf8]{inputenc}
\usepackage{lmodern}
\usepackage{textcomp}
\usepackage{enumitem}
\usepackage{geometry}
\geometry{margin=1in}
\begin{document}
\begin{center}
Answer Key - Geography - K-12 Level Assessment 6 \
\small (Answers are ordered exactly as in the question paper.)
\end{center}

\section*{Multiple Choice}
\begin{enumerate}[label=\arabic*.]
\item \textbf{Answer:} D
\\ \textit{Options:} A) Urban sprawl; B) Traffic increase; C) Traffic congestion; D) Decreased dependency on cars
\\ \textit{Note:} The construction of freeways typically leads to increased dependency on cars, as they facilitate faster and more convenient automobile travel, rather than decreasing it.
\item \textbf{Answer:} C
\\ \textit{Options:} A) Calcutta.; B) Bombay.; C) Cambodia.; D) Bali.
\\ \textit{Note:} The largest Hindu temple complex ever constructed, Angkor Wat, is located in Cambodia, showcasing the significance of Hindu architecture in Southeast Asia and the historical influence of Hinduism in the region.
\item \textbf{Answer:} B
\\ \textit{Options:} A) Ghettos; B) Suburbs; C) CBDs; D) Transition zones
\\ \textit{Note:} Planned communities are typically developed in suburban areas where land is more available for large-scale development and urban design, allowing for the creation of cohesive residential neighborhoods with specific amenities and infrastructure.
\item \textbf{Answer:} D
\\ \textit{Options:} A) Devolution; B) Colonialism; C) Heartland theory; D) Containment theory
\\ \textit{Note:} Containment theory was used as justification for US involvement in Vietnam because it aimed to prevent the spread of communism in Southeast Asia, reflecting the Cold War strategy to stop Soviet influence globally.
\item \textbf{Answer:} C
\\ \textit{Options:} A) manufacturing goods where production costs are the highest.; B) conducting accounting and research services where economical.; C) based on comparative advantage.; D) with headquarters located in LDCs.
\\ \textit{Note:} Transnational corporations distribute their operations based on comparative advantage to optimize efficiency and reduce costs by locating production and services in regions where they can be performed most effectively.
\end{enumerate}

\section*{True / False}
\begin{enumerate}[label=\arabic*.]
\setcounter{enumi}{5}
\item \textbf{Answer:} True
\\ \textit{Options:} A) True; B) False
\\ \textit{Note:} Deep dissection and erosion are more characteristic of youthful landscapes shaped by tectonic activity or steep gradients, while alluvial landscapes are typically formed by sediment deposition in flatter areas where water flows slowly.
\item \textbf{Answer:} True
\\ \textit{Options:} A) True; B) False
\\ \textit{Note:} National sports teams, flags, and anthems serve as powerful symbols that foster a sense of unity and pride among citizens, thereby promoting nationalism by representing a shared identity and collective values of the country.
\item \textbf{Answer:} True
\\ \textit{Options:} A) True; B) False
\\ \textit{Note:} The first zone in Burgess's concentric zone model, known as the Zone of Transition, is characterized by low-income housing, slums, and ethnic ghettos as it is typically the area where new immigrants settle and where urban decay occurs.
\item \textbf{Answer:} False
\\ \textit{Options:} A) True; B) False
\\ \textit{Note:} The number of frost-free days per year is crucial for agriculture because it directly affects the growing season, crop selection, and overall agricultural productivity in a country.
\item \textbf{Answer:} False
\\ \textit{Options:} A) True; B) False
\\ \textit{Note:} Young people are primarily drawn to major cities in Latin America for opportunities such as jobs, education, and lifestyle rather than solely because their families live there.
\end{enumerate}

\section*{Long Form Response}
\begin{enumerate}[label=\arabic*.]
\setcounter{enumi}{10}
\item \textbf{Answer:} Europe faces several major challenges today, including economic instability, political fragmentation, and the ongoing impacts of migration. These issues often demand urgent attention and action from governments and societies. In contrast, environmental apathy- where individuals and communities show indifference toward environmental issues-can hinder progress in addressing climate change and sustainability. This disinterest could have serious implications for the continent's future, as a lack of action on environmental concerns may exacerbate existing problems, leading to more severe economic and social consequences. Therefore, while Europe grapples with pressing challenges, overcoming environmental apathy is crucial for ensuring a sustainable and prosperous future.
\\ \textit{Note:} The answer correctly identifies Europe's significant challenges, such as economic instability and political fragmentation, while contrasting these immediate concerns with environmental apathy, emphasizing that indifference to environmental issues can impede progress on critical sustainability efforts essential for the continent's future.
\item \textbf{Answer:} Identifying perceptual or vernacular regions in geography presents several challenges, mainly due to the subjectivity involved in defining the criteria. Different researchers may have varying perspectives on what characteristics make up a region, such as cultural traits, economic activities, or even emotional connections to a place. This subjectivity can lead to inconsistencies in how regions are classified, making it difficult to reach a consensus among geographers. For example, one researcher might define a region based on local traditions, while another might focus on geographical features. As a result, the lack of a standardized approach complicates the study of these regions and can affect the accuracy of geographical analyses.
\\ \textit{Note:} correct because it emphasizes that the subjective nature of defining criteria for perceptual or vernacular regions results in varied interpretations among researchers, leading to inconsistencies in classification and challenges in achieving consensus.
\end{enumerate}

\vfill
\noindent\textit{End of Answer Key}
\end{document}